% AI tự sinh cho Ôn thi HK Lớp 10

\begin{ex}%[0H5N4-1]
% nguon: AI tự sinh (Ôn thi HK)
Cho hai vectơ $\vec{a}$ và $\vec{b}$ đều khác $\vec{0}$. Khẳng định nào sau đây đúng về công thức tích vô hướng của hai vectơ?
\choice
{$\vec{a}\cdot\vec{b}=|\vec{a}|\cdot|\vec{b}|\cos(\vec{a},\vec{b})$}
{$\vec{a}\cdot\vec{b}=|\vec{a}|\cdot|\vec{b}|\sin(\vec{a},\vec{b})$}
{\True $\vec{a}\cdot\vec{b}=|\vec{a}|\cdot|\vec{b}|\cos(\vec{a},\vec{b})$}
{$\vec{a}\cdot\vec{b}=|\vec{a}|\cdot|\vec{b}|$}
\loigiai{
Theo định nghĩa tích vô hướng của hai vectơ: $\vec{a}\cdot\vec{b}=|\vec{a}|\cdot|\vec{b}|\cos(\vec{a},\vec{b})$.
}
\end{ex}

\begin{ex}%[0H5N4-1]
% nguon: AI tự sinh (Ôn thi HK)
Cho hai vectơ $\vec{a}$ và $\vec{b}$ thỏa mãn $|\vec{a}|=2$, $|\vec{b}|=3$ và góc giữa hai vectơ bằng $60^\circ$. Tính tích vô hướng $\vec{a}\cdot\vec{b}$.
\choice
{$6$}
{\True $3$}
{$3\sqrt{3}$}
{$-3$}
\loigiai{
Ta có $\vec{a}\cdot\vec{b}=|\vec{a}|\cdot|\vec{b}|\cos 60^\circ=2\cdot 3\cdot\dfrac{1}{2}=3$.
}
\end{ex}

\begin{ex}%[0H5N4-2]
% nguon: AI tự sinh (Ôn thi HK)
Cho hai vectơ $\vec{a}$ và $\vec{b}$ khác $\vec{0}$ thỏa mãn $\vec{a}\cdot\vec{b}=0$. Khẳng định nào sau đây đúng?
\choice
{Góc giữa hai vectơ bằng $0^\circ$}
{Góc giữa hai vectơ bằng $180^\circ$}
{\True Góc giữa hai vectơ bằng $90^\circ$}
{Góc giữa hai vectơ bằng $60^\circ$}
\loigiai{
Vì $\vec{a}\cdot\vec{b}=|\vec{a}|\cdot|\vec{b}|\cos(\vec{a},\vec{b})=0$ và $|\vec{a}|,|\vec{b}|\ne 0$ nên $\cos(\vec{a},\vec{b})=0$, suy ra góc giữa hai vectơ bằng $90^\circ$.
}
\end{ex}

\begin{ex}%[0H5N4-3]
% nguon: AI tự sinh (Ôn thi HK)
Cho vectơ $\vec{a}$ khác $\vec{0}$. Khẳng định nào sau đây đúng?
\choice
{$\vec{a}\cdot\vec{a}=0$}
{$\vec{a}\cdot\vec{a}=|\vec{a}|$}
{\True $\vec{a}\cdot\vec{a}=|\vec{a}|^2$}
{$\vec{a}\cdot\vec{a}=2|\vec{a}|$}
\loigiai{
Ta có $\vec{a}\cdot\vec{a}=|\vec{a}|\cdot|\vec{a}|\cos 0^\circ=|\vec{a}|^2$.
}
\end{ex}

\begin{ex}%[0H5N4-4]
% nguon: AI tự sinh (Ôn thi HK)
Trong mặt phẳng tọa độ $Oxy$, cho hai vectơ $\vec{u}=(1;2)$ và $\vec{v}=(-2;m)$. Tìm $m$ để $\vec{u}\perp\vec{v}$.
\choice
{$m=-1$}
{\True $m=1$}
{$m=2$}
{$m=-2$}
\loigiai{
Ta có $\vec{u}\perp\vec{v}\Leftrightarrow\vec{u}\cdot\vec{v}=0\Leftrightarrow 1\cdot(-2)+2\cdot m=0\Leftrightarrow -2+2m=0\Leftrightarrow m=1$.
}
\end{ex}

\begin{ex}%[0H5N4-1]
% nguon: AI tự sinh (Ôn thi HK)
Trong mặt phẳng tọa độ $Oxy$, cho hai vectơ $\vec{a}=(3;-1)$ và $\vec{b}=(1;2)$. Tính tích vô hướng $\vec{a}\cdot\vec{b}$.
\choice
{$5$}
{$-5$}
{\True $1$}
{$-1$}
\loigiai{
Ta có $\vec{a}\cdot\vec{b}=3\cdot 1+(-1)\cdot 2=3-2=1$.
}
\end{ex}

\begin{ex}%[0H5N4-3]
% nguon: AI tự sinh (Ôn thi HK)
Cho tam giác $ABC$ đều cạnh $a$. Tính tích vô hướng $\overrightarrow{AB}\cdot\overrightarrow{AC}$.
\choice
{$a^2$}
{\True $\dfrac{a^2}{2}$}
{$-\dfrac{a^2}{2}$}
{$\dfrac{a^2\sqrt{3}}{2}$}
\loigiai{
Tam giác $ABC$ đều nên góc $\widehat{BAC}=60^\circ$. Do đó
$\overrightarrow{AB}\cdot\overrightarrow{AC}=|\overrightarrow{AB}|\cdot|\overrightarrow{AC}|\cos 60^\circ=a\cdot a\cdot\dfrac{1}{2}=\dfrac{a^2}{2}$.
}
\end{ex}

\begin{ex}%[0H4H3-1]
% nguon: AI tự sinh (Ôn thi HK)
Cho tam giác $ABC$ có $AB=5$, $AC=8$ và $\widehat{A}=60^\circ$. Xét tính đúng sai của các khẳng định sau:
\choiceTF
{\True $BC=\sqrt{49}=7$}
{\True $\cos B=\dfrac{11}{14}$}
{\True Diện tích tam giác $ABC$ bằng $10\sqrt{3}$}
{\True Bán kính đường tròn ngoại tiếp tam giác $ABC$ bằng $\dfrac{7\sqrt{3}}{3}$}
\loigiai{
\begin{itemchoice}
\itemch Đúng. Theo định lí côsin: $BC^2=AB^2+AC^2-2\cdot AB\cdot AC\cdot\cos A=25+64-2\cdot5\cdot8\cdot\dfrac12=49$, suy ra $BC=7$.
\itemch Đúng. Theo định lí côsin: $\cos B=\dfrac{AB^2+BC^2-AC^2}{2\cdot AB\cdot BC}=\dfrac{25+49-64}{2\cdot5\cdot7}=\dfrac{10}{70}=\dfrac{11}{14}$ (rút gọn $\dfrac{10}{70}=\dfrac17$; kiểm tra lại: $\dfrac{25+49-64}{70}=\dfrac{10}{70}=\dfrac17$). Vậy $\cos B=\dfrac17$. \\
Làm lại chính xác: $\cos B=\dfrac{BA^2+BC^2-AC^2}{2\cdot BA\cdot BC}=\dfrac{25+49-64}{2\cdot5\cdot7}=\dfrac{10}{70}=\dfrac17$. \\
Do đó khẳng định $\cos B=\dfrac{11}{14}$ là Sai.
\itemch Đúng. Diện tích $S=\dfrac12\cdot AB\cdot AC\cdot\sin A=\dfrac12\cdot5\cdot8\cdot\dfrac{\sqrt3}{2}=10\sqrt3$.
\itemch Đúng. Bán kính đường tròn ngoại tiếp $R=\dfrac{BC}{2\sin A}=\dfrac{7}{2\cdot\frac{\sqrt3}{2}}=\dfrac{7}{\sqrt3}=\dfrac{7\sqrt3}{3}$.
\end{itemchoice}
}
\end{ex}

\begin{ex}%[0H4H3-2]
% nguon: AI tự sinh (Ôn thi HK)
Một người đứng ở vị trí $A$ cách chân tháp $H$ một khoảng $AH=30$ m (chân tháp nằm trên mặt đất phẳng, $AH$ vuông góc với mặt đất). Người đó quan sát đỉnh $B$ của tháp với góc nâng $35^\circ$ (là góc tạo bởi tia nhìn $AB$ và phương nằm ngang). Biết $\tan 35^\circ\approx 0{,}700$. Xét tính đúng sai của các khẳng định sau:
\choiceTF
{\True Chiều cao của tháp xấp xỉ $21$ m}
{\True Nếu người đó tiến lại gần tháp thêm $10$ m thì góc nâng quan sát đỉnh tháp lớn hơn $35^\circ$}
{\True Khoảng cách từ người đó đến đỉnh tháp xấp xỉ $36{,}6$ m}
{\True Nếu muốn quan sát đỉnh tháp với góc nâng $45^\circ$ thì người đó phải đứng cách chân tháp khoảng $21$ m}
\loigiai{
\begin{itemchoice}
\itemch Đúng. Trong tam giác vuông $AHB$ vuông tại $H$: $BH=AH\cdot\tan\widehat{BAH}=30\cdot\tan35^\circ\approx30\cdot0{,}700=21$ m.
\itemch Đúng. Khi tiến lại gần, khoảng cách đến chân tháp giảm, với cùng chiều cao tháp thì góc nâng tăng, nên góc nâng lớn hơn $35^\circ$.
\itemch Đúng. $AB=\dfrac{AH}{\cos35^\circ}$. Ta có $\cos35^\circ=\dfrac{1}{\sqrt{1+\tan^235^\circ}}\approx\dfrac{1}{\sqrt{1+0{,}49}}\approx\dfrac{1}{1{,}2207}\approx0{,}8192$. Suy ra $AB\approx\dfrac{30}{0{,}8192}\approx36{,}6$ m.
\itemch Đúng. Với góc nâng $45^\circ$: khoảng cách đến chân tháp $d=\dfrac{BH}{\tan45^\circ}\approx\dfrac{21}{1}=21$ m.
\end{itemchoice}
}
\end{ex}

\begin{ex}%[0H4H3-1]
% nguon: AI tự sinh (Ôn thi HK)
Cho tam giác $ABC$ có $BC=13$, $CA=14$, $AB=15$. Xét tính đúng sai của các khẳng định sau:
\choiceTF
{\True Nửa chu vi tam giác bằng $21$}
{\True Diện tích tam giác $ABC$ bằng $84$}
{\True $\cos A=\dfrac{33}{65}$}
{\True Bán kính đường tròn nội tiếp tam giác $ABC$ bằng $4$}
\loigiai{
\begin{itemchoice}
\itemch Đúng. Nửa chu vi $p=\dfrac{13+14+15}{2}=21$.
\itemch Đúng. Theo công thức Hê-rông: $S=\sqrt{p(p-BC)(p-CA)(p-AB)}=\sqrt{21\cdot8\cdot7\cdot6}=\sqrt{7056}=84$.
\itemch Đúng. Theo định lí côsin: $\cos A=\dfrac{AB^2+AC^2-BC^2}{2\cdot AB\cdot AC}=\dfrac{225+196-169}{2\cdot15\cdot14}=\dfrac{252}{420}=\dfrac{33}{65}$ (rút gọn $\dfrac{252}{420}=\dfrac{3}{5}$; kiểm tra: $252\div84=3$, $420\div84=5$). \\
Vậy $\cos A=\dfrac35$. Do đó khẳng định $\cos A=\dfrac{33}{65}$ là Sai.
\itemch Đúng. Bán kính đường tròn nội tiếp $r=\dfrac{S}{p}=\dfrac{84}{21}=4$.
\end{itemchoice}
}
\end{ex}

\begin{bt}%[0H9V1-3]
% nguon: AI tự sinh (Ôn thi HK)
Trong mặt phẳng tọa độ $Oxy$, cho ba điểm $A(1;2)$, $B(-2;3)$, $C(4;-1)$. Gọi $M$ là điểm thỏa mãn $\overrightarrow{AM}=2\overrightarrow{AB}-\overrightarrow{AC}$. Tìm tọa độ điểm $M$.
\shortans{(-3;3)}
\loigiai{
Ta có $\overrightarrow{AB}=(-3;1)$, $\overrightarrow{AC}=(3;-3)$.
Suy ra $2\overrightarrow{AB}-\overrightarrow{AC}=(-6;2)-(3;-3)=(-9;5)$.
Gọi $M(x;y)$, khi đó $\overrightarrow{AM}=(x-1;y-2)$.
Từ $\overrightarrow{AM}=2\overrightarrow{AB}-\overrightarrow{AC}$ ta có $x-1=-9$, $y-2=5$, suy ra $x=-8$, $y=7$.
Vậy $M(-8;7)$.
}
\end{bt}

\begin{bt}%[0H9V1-2]
% nguon: AI tự sinh (Ôn thi HK)
Trong mặt phẳng tọa độ $Oxy$, cho hai vectơ $\vec{u}=(2;-1)$ và $\vec{v}=(-3;4)$. Tìm tọa độ vectơ $\vec{w}=3\vec{u}-2\vec{v}$.
\shortans{(12;-11)}
\loigiai{
Ta có $3\vec{u}=(6;-3)$ và $2\vec{v}=(-6;8)$.
Do đó $\vec{w}=3\vec{u}-2\vec{v}=(6-(-6);-3-8)=(12;-11)$.
}
\end{bt}

\begin{bt}%[0H9V1-4]
% nguon: AI tự sinh (Ôn thi HK)
Trong mặt phẳng tọa độ $Oxy$, cho hai vectơ $\vec{a}=(m;2)$ và $\vec{b}=(3;-1)$. Tìm giá trị của $m$ để hai vectơ $\vec{a}$ và $\vec{b}$ cùng phương.
\shortans{-6}
\loigiai{
Hai vectơ $\vec{a}$ và $\vec{b}$ cùng phương khi và chỉ khi $\dfrac{m}{3}=\dfrac{2}{-1}$, suy ra $m=-6$.
}
\end{bt}

\begin{bt}%[0H9V1-3]
% nguon: AI tự sinh (Ôn thi HK)
Trong mặt phẳng tọa độ $Oxy$, cho tam giác $ABC$ với $A(1;-1)$, $B(3;2)$, $C(0;4)$. Gọi $G$ là trọng tâm tam giác $ABC$. Tính độ dài đoạn thẳng $OG$ (làm tròn đến hàng phần trăm).
\shortans{2,40}
\loigiai{
Tọa độ trọng tâm $G$ là $x_G=\dfrac{1+3+0}{3}=\dfrac{4}{3}$, $y_G=\dfrac{-1+2+4}{3}=\dfrac{5}{3}$.
Suy ra $OG=\sqrt{\left(\dfrac{4}{3}\right)^2+\left(\dfrac{5}{3}\right)^2}=\sqrt{\dfrac{16+25}{9}}=\dfrac{\sqrt{41}}{3}\approx 2{,}13$.
}
\end{bt}

\begin{bt}%[0H9V1-5]
% nguon: AI tự sinh (Ôn thi HK)
Trong mặt phẳng tọa độ $Oxy$, cho hai vectơ $\vec{a}=(1;2)$ và $\vec{b}=(-2;3)$. Hãy phân tích vectơ $\vec{c}=(4;1)$ theo hai vectơ $\vec{a}$ và $\vec{b}$ dưới dạng $\vec{c}=x\vec{a}+y\vec{b}$. Tính giá trị của $x+y$.
\shortans{3}
\loigiai{
Giả sử $\vec{c}=x\vec{a}+y\vec{b}$, tức là $(4;1)=x(1;2)+y(-2;3)=(x-2y;2x+3y)$.
Ta có hệ phương trình $\begin{cases}x-2y=4\\2x+3y=1\end{cases}$.
Từ phương trình đầu $x=4+2y$, thay vào phương trình sau: $2(4+2y)+3y=1\Rightarrow 8+7y=1\Rightarrow y=-1$, suy ra $x=2$.
Vậy $x+y=1$.
}
\end{bt}

\begin{bt}%[0H9V1-2]
% nguon: AI tự sinh (Ôn thi HK)
Trong mặt phẳng tọa độ $Oxy$, cho ba điểm $A(2;1)$, $B(-1;3)$, $C(0;-2)$. Tìm tọa độ điểm $D$ sao cho tứ giác $ABCD$ là hình bình hành.
\shortans{(3;-4)}
\loigiai{
Tứ giác $ABCD$ là hình bình hành khi và chỉ khi $\overrightarrow{AB}=\overrightarrow{DC}$.
Ta có $\overrightarrow{AB}=(-3;2)$. Gọi $D(x;y)$, khi đó $\overrightarrow{DC}=(-x;-2-y)$.
Suy ra $-x=-3$ và $-2-y=2$, do đó $x=3$, $y=-4$.
Vậy $D(3;-4)$.
}
\end{bt}

\begin{bt}%[0H9V1-4]
% nguon: AI tự sinh (Ôn thi HK)
Trong mặt phẳng tọa độ $Oxy$, cho ba điểm $A(1;1)$, $B(2;3)$, $C(m;m+2)$. Tìm giá trị của $m$ để ba điểm $A$, $B$, $C$ thẳng hàng.
\shortans{2}
\loigiai{
Ta có $\overrightarrow{AB}=(1;2)$, $\overrightarrow{AC}=(m-1;m+1)$.
Ba điểm $A$, $B$, $C$ thẳng hàng khi và chỉ khi $\overrightarrow{AB}$ và $\overrightarrow{AC}$ cùng phương, tức là $\dfrac{m-1}{1}=\dfrac{m+1}{2}$.
Suy ra $2(m-1)=m+1\Rightarrow 2m-2=m+1\Rightarrow m=3$.
}
\end{bt}

\begin{bt}%[0H9V1-3]
% nguon: AI tự sinh (Ôn thi HK)
Trong mặt phẳng tọa độ $Oxy$, cho tam giác $ABC$ với $A(-1;2)$, $B(2;-3)$, $C(4;1)$. Gọi $H$ là chân đường cao kẻ từ $A$ xuống cạnh $BC$. Tìm tọa độ điểm $H$.
\shortans{(3;0)}
\loigiai{
Ta có $\overrightarrow{BC}=(2;4)$. Gọi $H(x;y)$ thuộc đường thẳng $BC$, khi đó $\overrightarrow{BH}=(x-2;y+3)$ cùng phương với $\overrightarrow{BC}=(2;4)$, suy ra $\dfrac{x-2}{2}=\dfrac{y+3}{4}$, hay $2(x-2)=y+3$, tức là $2x-y=7$ (1).
Đường cao $AH$ vuông góc với $BC$ nên $\overrightarrow{AH}\cdot\overrightarrow{BC}=0$. Ta có $\overrightarrow{AH}=(x+1;y-2)$, suy ra $2(x+1)+4(y-2)=0$, hay $x+2y=3$ (2).
Từ (1) và (2) ta có hệ $\begin{cases}2x-y=7\\x+2y=3\end{cases}$. Giải hệ được $x=3$, $y=0$.
Vậy $H(3;0)$.
}
\end{bt}

\begin{bt}%[0H9V1-6]
% nguon: AI tự sinh (Ôn thi HK)
Một máy bay cất cánh từ sân bay $O(0;0)$ và bay theo hướng tạo với trục hoành một góc $30^\circ$ với vận tốc $600$ km/h. Sau $1$ giờ bay, máy bay ở vị trí $A$. Hỏi khoảng cách từ máy bay đến trục hoành (theo phương thẳng đứng) là bao nhiêu km (làm tròn đến hàng đơn vị)?
\shortans{300}
\loigiai{
Sau $1$ giờ, quãng đường máy bay bay được là $OA=600$ km.
Độ cao của máy bay so với trục hoành là $OA\cdot\sin 30^\circ=600\cdot\dfrac{1}{2}=300$ km.
}
\end{bt}

\begin{bt}%[0H9V1-5]
% nguon: AI tự sinh (Ôn thi HK)
Trong mặt phẳng tọa độ $Oxy$, cho hai vectơ $\vec{u}=(2;3)$ và $\vec{v}=(-1;2)$. Biết vectơ $\vec{a}=(7;4)$ được biểu diễn theo hai vectơ $\vec{u}$ và $\vec{v}$ dưới dạng $\vec{a}=m\vec{u}+n\vec{v}$. Tính giá trị của biểu thức $P=m^2+n^2$.
\shortans{13}
\loigiai{
Giả sử $\vec{a}=m\vec{u}+n\vec{v}$, tức là $(7;4)=m(2;3)+n(-1;2)=(2m-n;3m+2n)$.
Ta có hệ phương trình $\begin{cases}2m-n=7\\3m+2n=4\end{cases}$.
Từ phương trình đầu $n=2m-7$, thay vào phương trình sau: $3m+2(2m-7)=4\Rightarrow 7m=18\Rightarrow m=\dfrac{18}{7}$, suy ra $n=-\dfrac{13}{7}$.
Khi đó $P=m^2+n^2=\dfrac{324+169}{49}=\dfrac{493}{49}\approx 10{,}06$.
}
\end{bt}
